\section*{Capítulo \ref{ch:b3:green-industrial} --- Química verde y procesos industriales}

\begin{solution}{exo:b3:green-industrial:1}
Diels--Alder: 100\,\% (una adición). Esterificación: $88.0/(60.0 + 46.0) = 83$\,\%. Wittig:
$104.0/(106.0 + 276.0) = 27$\,\%; el óxido de trifenilfosfina se lleva la mayor parte de la masa.
\end{solution}

\begin{solution}{exo:b3:green-industrial:2}
$\mathrm{PMI} = 120/8.0 = 15$; $E = \mathrm{PMI} - 1 = 14$.
\end{solution}

\begin{solution}{exo:b3:green-industrial:3}
Un lavado de una carga normal a una temperatura y con un resultado de limpieza dados. Los detergentes
se dosifican de forma distinta: un polvo concentrado necesita menos por lavado, de modo que masas iguales
no prestan el mismo servicio.
\end{solution}

\begin{solution}{exo:b3:green-industrial:4}
2-Metiltetrahidrofurano o acetato de etilo en lugar de diclorometano; heptano (con acetato
de etilo) en lugar de hexano; en lugar de la DMF, acetato de etilo, 2-metiltetrahidrofurano o
carbonato de dimetilo, o un acoplamiento en agua con \omterm{def:b3:colloids:surfactant}{tensioactivos}, según los
reactivos.
\end{solution}

\begin{solution}{exo:b3:green-industrial:5}
1.000 mol de ácido; 0.800 mol de éster (\qty{88.0}{g/mol}): rendimiento del 80\,\%. AE: 83.0\,\%. SF $= 152/106 = 1.43$.
$\mathrm{RME} = 70.4/152 = 0.463$; $\mathrm{AE} \times \text{rendimiento}/\mathrm{SF} = 0.830 \times 0.80/1.43 = 0.463$.
\end{solution}

\begin{solution}{exo:b3:green-industrial:6}
Clorhidrina: $44.0/(28.0 + 71.0 + 74.1) = 25$\,\%; oxidación directa: 100\,\%. Un mol de $\ce{CaCl2}$ por mol de
óxido de etileno: $\qty{1.0e6}{g}/\qty{44.0}{g/mol} \times \qty{111.1}{g/mol} = \qty{2.5}{t}$ por tonelada.
\end{solution}

\begin{solution}{exo:b3:green-industrial:7}
$3/8 \times \qty{5.88e4}{mol} \times \qty{44.0}{g/mol} = \qty{0.97}{t}$ de $\ce{CO2}$ por tonelada de amoníaco. Para el hidrógeno, $\ce{CH4 + 2H2O -> CO2 + 4H2}$:
$1/4$ mol por mol, $\qty{5.0e5}{mol}/4 \times \qty{44.0}{g/mol} = \qty{5.5}{t}$ de $\ce{CO2}$ por tonelada de hidrógeno.
\end{solution}

\begin{solution}{exo:b3:green-industrial:8}
% ledger: dfg:H2O_l, dfh:H2O-l
Un kilogramo son \qty{500}{mol}. A partir de $\Delta_rG^\circ = \qty{237.1}{kJ/mol}$: $\qty{1.19e5}{kJ} = \qty{32.9}{kWh}$. A partir de $\Delta_rH^\circ = \qty{285.8}{kJ/mol}$: \qty{39.7}{kWh}
(la diferencia es calor que puede aportar el entorno).
\end{solution}

\begin{solution}{exo:b3:green-industrial:9}
$\qty{1.0e6}{g}/\qty{146.0}{g/mol} = \qty{6.85e3}{mol}$ de $\ce{N2O}$, $\qty{0.30}{t}$; $\times 273 = \qty{82}{t}$ de $\ce{CO2}$ equivalente por tonelada de ácido adípico.
\end{solution}

\begin{solution}{exo:b3:green-industrial:10}
La RME de una etapa cuenta como consumido el reactivo en exceso; si el exceso se recupera
y se realimenta, solo se consume la parte realmente perdida, y la eficiencia másica global
es mayor que el producto de los valores de las etapas. En el \omterm{def:b3:green-industrial:e-factor}{factor E}, la
corriente reciclada no es un residuo, siempre que el reciclado esté dentro del límite;
hay que contabilizar entonces su energía y sus pérdidas.
\end{solution}

\begin{solution}{exo:b3:green-industrial:11}
Antes: $20 \times 0.80 = \qty{16}{kg}$ de disolvente, con un 10\,\% perdido: \qty{1.6}{kg} de residuo. Después: \qty{4.0}{kg}, con \qty{0.40}{kg} perdidos. El
\omterm{def:b3:green-industrial:e-factor}{factor E} completo disminuye en 1.2 kg por kilogramo de producto (y la energía de destilación,
en tres cuartas partes).
\end{solution}

\begin{solution}{exo:b3:green-industrial:12}
Por \omterm{def:b3:green-industrial:lca}{unidad funcional} y con el mismo límite, presentaría cada categoría de impacto
por separado (clima, uso del suelo, uso del agua\dots) en lugar de una sola puntuación. Una
decisión necesita los pesos atribuidos a esas categorías, la escasez local de agua
y de suelo, la incertidumbre de los datos y saber si la materia prima de base biológica
compite con la alimentación.
\end{solution}

\begin{solution}{pb:b3:green-industrial:1}
\textbf{1.} $\ce{C10H14 + C4H6O3 -> C12H16O + C2H4O2}$ (HF); $\ce{C12H16O + H2 -> C12H18O}$ (níquel Raney);
$\ce{C12H18O + CO -> C13H18O2}$ (paladio).

\textbf{2.} $\ce{C10H14 + C4H6O3 + H2 + CO -> C13H18O2 + C2H4O2}$.

\textbf{3.} $134.0 + 102.0 + 2.0 + 28.0 = \qty{266.0}{g/mol}$ de reactivos; ibuprofeno, \qty{206.0}{g/mol}: AE $= 206.0/266.0 = 77$\,\%.

\textbf{4.} Isobutilbenceno, anhídrido acético, cloroacetato de etilo, etóxido de sodio,
ácido acuoso ($\ce{H3O+}$), hidroxilamina y agua.

\textbf{5.} $134.0 + 102.0 + 122.5 + 68.0 + 19.0 + 33.0 + 2 \times 18.0 = \qty{514.5}{g/mol}$: AE $= 206.0/514.5 = 40$\,\%.

\textbf{6.} Ácido acético, cloruro de sodio, etanol, dióxido de carbono, agua, amoníaco (en forma de
sal de amonio) y sales de aluminio procedentes del cloruro de aluminio estequiométrico.

\textbf{7.} $0.71 + 0.55 + 0.010 + 0.14 + 0.05 = \qty{1.46}{t}$.

\textbf{8.} PMI: 1.46; $E = 0.46$.

\textbf{9.} $0.46 - 0.31 = 0.15$.

\textbf{10.} PMI $= 2.0 + 1.5 = 3.5$; $E = 2.5$.

\textbf{11.} Los rendimientos inferiores al 100\,\%, los reactivos en exceso, las pérdidas de disolvente y de catalizador y los
materiales de tratamiento añaden residuos que la \omterm{def:b3:green-industrial:atom-economy}{economía atómica} ignora.

\textbf{12.} Prevenir los residuos, maximizar la \omterm{def:b3:green-industrial:atom-economy}{economía atómica}, usar catalizadores en lugar de
reactivos estequiométricos, evitar derivados y etapas innecesarios, ahorrar energía.

\textbf{13.} El cloruro de aluminio estequiométrico, hidrolizado a residuos de aluminio en el
tratamiento; el fluoruro de hidrógeno es a la vez catalizador y disolvente, y se recupera y se reutiliza.

\textbf{14.} El níquel Raney (un catalizador de hidrogenación heterogéneo).

\textbf{15.} Un complejo fosfina--paladio: la química de \omterm{def:b3:homogeneous-catalysis:carbonylation}{carbonilación} del
\cref{ch:b3:homogeneous-catalysis}.

\textbf{16.} El anhídrido acético transfiere un grupo acetilo; la otra mitad de la
molécula sale como ácido acético.

\textbf{17.} El fluoruro de hidrógeno es muy tóxico y corrosivo, y el monóxido de carbono, tóxico e
inflamable: equipos cerrados de materiales resistentes, detectores y operarios
formados.

\textbf{18.} Cada etapa añade calentamiento, enfriamiento, separaciones y recuperación de disolventes; tres
etapas necesitan menos de todo ello.

\textbf{19.} $2.5 \times 3000 = \qty{7500}{t}$ al año.

\textbf{20.} $0.15 \times 3000 = \qty{450}{t}$ al año.

\textbf{21.} Unas \qty{7000}{t} de residuos evitados cada año.

\textbf{22.} No: importa la naturaleza de los residuos (una tonelada de agua salada no es una tonelada de
un compuesto tóxico persistente), y también la energía, las emisiones y los peligros de los
reactivos.

\textbf{23.} Los impactos de la fabricación aguas arriba de los reactivos y de la energía, de las emisiones al
aire y al agua, y del uso y la eliminación del producto, por \omterm{def:b3:green-industrial:lca}{unidad funcional}, en
varias categorías de impacto.

\textbf{24.} La ruta catalítica en tres etapas tiene una \omterm{def:b3:green-industrial:atom-economy}{economía atómica} de alrededor del 77\,\% (40\,\% para la
ruta en seis etapas).
\end{solution}
